% 一、问题重述
\section{问题重述}
\subsection{问题背景}
非孤岛式微网由小区负载、光伏发电、储能设备和外部电网共同构成。光伏出力具有明显的时变性，储能设备可以在电价较低或系统存在富余电量时充电，并在高价或供能不足时释放电量，从而在满足负载需求的前提下降低外网购电费用。微网调控需要根据负载、光伏出力、储能状态以及可获得的电价和预报信息，合理确定各时段购电量与储能充放电量。

题设储能最大容量为 12000\,kWh，运行储电量限制在 1200--10800\,kWh，最大充放电功率为 5000\,kW，2025 年 1 月 1 日 0:00 储电量为 6000\,kWh，充放电效率均为 90\%。附件 1 给出典型日的电价、负载及光伏预测功率；附件 2 给出 2025 年全年负载与光伏实际功率；附件 3 给出每天 0:00、6:00、12:00、18:00 发布的未来 24\,h 整点光伏预报；附件 4 给出全年实时电价。随着可用信息和交易权限增加，题目依次要求建立四类购电与储能调度策略。

\subsection{问题要求}
\textbf{问题 1}\quad 每日负载和电价保持一致，并给定一天的光伏预测曲线。要求每天 0:00 一次性确定全天计划购电策略，在各时段微网供能不低于负载、储能满足容量与功率限制且 0:00 与 24:00 储电量相同的条件下，使全天计划购电费用尽可能低，并给出指定时段购电量、全天购电量与购电费用，以及各指定区间的储能充放电量。

\textbf{问题 2}\quad 电价仍采用附件 1 给定的日内曲线，负载和光伏出力随日期变化。每天 0:00 确定全天计划购电量，计划确定后不再调整；实际运行中若微网供能出现缺口，则按交易时刻电价的 5 倍向外网紧急购电。需要结合附件 2 构造逐日购电策略，计算指定日期的计划购电、充放电和紧急购电结果，并保存 2 月 1 日至 12 月 31 日的完整策略。

\textbf{问题 3}\quad 在问题 2 基础上，每天 0:00、6:00、12:00 和 18:00 均可获得未来 24\,h 整点光伏预报。0:00 预报用于制定全天计划，后续预报可用于调整尚未执行的购电量。调减部分按交易时刻电价的 50\% 支付违约费用，超过原计划的调增部分按 1.5 倍交易时刻电价购入，总费用还包括正常计划购电和紧急购电费用。需要建立能够随预报更新进行调整的购电策略，并定量判断引入其他时刻预报的实际价值。

\textbf{问题 4}\quad 将外网电价进一步扩展为实时波动电价。结合附件 2 的实际负载与光伏数据、附件 3 的光伏预报以及附件 4 的全年电价，在波动电价条件下分别按照问题 2 和问题 3 的交易权限重新建立购电策略并计算相应结果。
